\documentclass[11pt]{article}

\usepackage[a4paper,margin=1in]{geometry}

\usepackage[T1]{fontenc}
\usepackage[utf8]{inputenc}

\usepackage{amsmath,amssymb,amsthm}
\usepackage{mathtools}

\usepackage{graphicx}
\usepackage{xcolor}
\usepackage{hyperref}

\usepackage{enumitem}
\usepackage{booktabs}
\usepackage{array}

\hypersetup{
    colorlinks=true,
    linkcolor=blue,
    urlcolor=blue,
    citecolor=blue
}

\title{Randomised Algorithms}
\author{Sahil M Deo}
\date{\today}

\begin{document}

\maketitle

\tableofcontents
\newpage

\section{Introduction}

\subsection{What even is a Randomised Algorithm?}

A randomised algorithm is an algorithm that makes random choices during its execution. Unlike deterministic algorithms, whose behaviour is completely determined by the input, the behaviour of a randomised algorithm may vary between different executions on the same input.

Typically, the algorithm uses a random number generator (rng for short) to decide between multiple possible actions. As a result, quantities such as running time, memory usage, or even correctness\footnote{
The fact that correctness itself may become probabilistic can initially seem undesirable, especially in critical sumYstems where absolute correctness is expected. However, in many practical situations, a sufficiently small error probability is considered acceptable. For example, if an algorithm can be shown to fail with probability at most \(10^{-30}\), then the chance of failure may already be smaller than the probability of random hardware faults such as memory bit flips caused by cosmic radiation, which programmers typically ignore during ordinary software development.
} can become random variables (though usually not all of them at the same time).


The performance of a randomised algorithm is therefore analysed using probability and expectation. If $T_{run}$ denotes the runtime of the algorithm, we are usually interested in quantities such as $T_{algo}=E[T_{run}]$ over all possible random choices made by the algorithm.\footnote{$T_{algo}$ can still be a random variable over the input - so this may \textbf{not} be the same as average case performance of the algorithm $T_{avg}$, which is an expected value over all possible inputs.}

\vspace{0.6em}
Randomised algorithms are often used for the following purposes:



\begin{itemize}
    \item avoid adversarial worst case inputs,
    \item simplify algorithm design,
    \item obtain fast approximate solutions.
\end{itemize}

\subsection{Implementation of Randomness in \texttt{C++}}

Lets make things more concrete with some actual code and implementation details!
Modern \texttt{C++} provides pseudo-random\footnote{Since computers are deterministic, they rely on parameters like atmospheric noise and temperature for randomness. The details are irrelevant for us... We will take pseudo random as true random for our purposes} number generators through the \texttt{<random>} library. A commonly used generator is \texttt{mt19937}, which is based on the Mersenne Twister algorithm.

\begin{verbatim}
#include <bits/stdc++.h>
mt19937 rng(chrono::steady_clock::now().time_since_epoch().count()); //using time as a seed
\end{verbatim}

This creates an object rng with seed\footnote{seed values are used in pseudo random generators as a way to recreate the generated random numbers. This can be used for debugging a randomised algorithm.} as current time which will be passed as a parameter to other functions. It has the capability to produce "truly" random 32-bit integers.

\subsubsection*{Generating Random Integers}

To generate a random integer uniformly from a range \([L,R]\), we use

\begin{verbatim}
uniform_int_distribution<int>uid(L,R);

int x=uid(rng);
int y=uid(rng);
.
.
.
\end{verbatim}

Each generation takes around constant time.

\subsubsection*{Generating Random Real Numbers}

To generate a real number uniformly from an interval \([x,y)\), we use

\begin{verbatim}
uniform_real_distribution<double>urd(x,y);

double x=urd(rng);
double y=urd(rng);
.
.
.
\end{verbatim}

This also takes \(\Theta(1)\) time per generation.

\subsubsection*{Generating a Random Permutation}

A random permutation of an array (in-place modification) can be generated using \texttt{shuffle}:

\begin{verbatim}
shuffle(a.begin(),a.end(),rng);
\end{verbatim}

Internally, this does something like:

\begin{verbatim}
for(int i=n-1;i>0;i--)
{
    int j=random integer from [0,i];
    swap(a[i],a[j]);
}
\end{verbatim}

(Try to think why all permutations are equally likely with this method!)

\subsubsection*{Implementing Random Decisions}

Sometimes, an algorithm must choose between multiple actions with specified probabilities. For example, let three actions X,Y,Z be such that they are executed with probabilities:

\[
\Pr(X)=\frac{1}{5},
\qquad
\Pr(Y)=\frac{3}{5},
\qquad
\Pr(Z)=\frac{1}{5}
\]

One simple way to implement this is to generate a random integer uniformly from \(1\) to \(5\):

\begin{verbatim}
uniform_int_distribution<int>uid(1,5);

int v=uid(rng);

if(v==1)
{
    //do X
}
else if(v<=4)
{
    //do Y
}
else
{
    //do Z
}
\end{verbatim}

thus implementing the required probabilities.

\section{Algorithmic Performance} 

For algorithms whose performance\footnote{
Although this section focuses primarily on running time, the same methods of analysis can also be applied to other resources such as memory usage.
} depends on the input size or other input parameters, we usually estimate the number of elementary operations performed as a function of the input size \(n\) or the maximum value of input $A_{max}$, usually denoted by $A$. This gives a mathematical way to compare algorithms independent of hardware or implementation details.
\\\\
A classical example is Bubble Sort. The algorithm repeatedly traverses the array, comparing adjacent elements and swapping them whenever they are in the wrong order. After each pass, the largest remaining element moves to its correct position near the end of the array.
\\\\
With some thought it can be proved that the total number of swaps performed by Bubble Sort is exactly equal to the number of inversions $N_{inv}$ in the input array. (An inversion is a pair of indices \((i,j)\) such that \(i<j\) but \(a_i>a_j\).) By taking swaps to be one operation taking constant time, we can estimate runtime of Bubble sort to be $T_{run}\propto N_{inv}$

\subsection{Worst Case Analysis}

In worst case analysis, we determine the maximum possible number of operations over all valid inputs of size \(n\). For Bubble Sort, the worst case occurs when the array is sorted in reverse order, where every pair is inverted. 

\[N_{inv,max}=\binom{n}{2}=
\frac{n(n-1)}{2}
\]

Worst case analysis is useful because it guarantees an upper bound on the running time regardless of the input provided.

\subsection{Average Case Analysis} 

In average case analysis, we compute the expected number of operations assuming that all valid inputs are equally likely.

For Bubble Sort, this corresponds to analysing the expected number of inversions in a random permutation of size \(n\). For every pair of indices \((i,j)\) with \(i<j\), the probability that the pair forms an inversion is

\[
\Pr\big((i,j)\text{ is an inversion}\big)=\frac{1}{2}
\]

\noindent Define the indicator variable:

\[
I_{ij}=
\begin{cases}
1,&\text{if }(i,j)\text{ is an inversion}\\
0,&\text{otherwise}
\end{cases}
\]

\[
N_{\mathrm{inv}}=\sum_{1\le i<j\le n}I_{ij}
\]

\[
\mathbb{E}[N_{\mathrm{inv}}]
=
\mathbb{E}\left[\sum_{1\le i<j\le n}I_{ij}\right]
=
\sum_{1\le i<j\le n}\mathbb{E}[I_{ij}]
\]
\[
\mathbb{E}[I_{ij}]
=
1\cdot\frac{1}{2}+0\cdot\frac{1}{2}
=
\frac{1}{2}
\]

Therefore,

\[
\mathbb{E}[N_{\mathrm{inv}}]
=
\sum_{1\le i<j\le n}\frac{1}{2}
=
\frac{1}{2}\binom{n}{2}
=
\frac{n(n-1)}{4}
\]

So notice that if we randomise the input sequence, $T_{algo}$ becomes equal to $T_{avg}$ thus improving performance by a factor\footnote{Randomising the input itself takes about n operations, so more rigorously we would say an asumYmptotic factor of 2. In most cases we will be able to neglect the time required to randomise the input.} of 2.
\subsection{Time Complexity Analysis}

We are sometimes interested in comparing runtimes like $T_{run}=n^2$ to $T_{run}=100nlogn$.
To compare growth rates, we use asumYmptotic notation:

\[
T(n)=\Theta(f(n))
\]

means that \(T(n)\) grows asumYmptotically at the same rate as \(f(n)\).

My favourite way to think about time complexity is as follows: we try to find the ``simplest'' function \(f(n)\) such that

\[
\lim_{n\to\infty}\frac{T(n)}{f(n)}=k,
\qquad
k\in\mathbb{R}_{>0}
\]

That is, \(T(n)\) and \(f(n)\) grow at the same asumYmptotic rate up to a positive constant factor.

For example, Bubble Sort performs on the order of \(n^2\) operations in both the average and worst case, so its time complexity is

\[
T(n)=\Theta(n^2)
\]

AsumYmptotic analysis allows us to compare algorithms by their rate of growth rather than implementation-dependent constants.

This allows us to say that A: $T_{run}=n^2$ is worse than B: $T_{run}=100nlogn$ since A has $T(n)=\Theta(n^2)$ and B has $T(n)=\Theta(nlogn)$ even though for some small n the first one may well be better, but for small n the performance doesn't really matter\dots

\section{Standard Examples}
\subsection{Quick Sort}

Quick Sort is one of the most important examples of a randomised algorithm in practice. The basic idea is very simple:

\begin{enumerate}
    \item Choose a pivot element.
    \item Partition the array into:
    \begin{itemize}
        \item elements smaller than the pivot,
        \item elements equal to the pivot,
        \item elements greater than the pivot.
    \end{itemize}
    This takes n comparisions with our pivot 
    \item Recursively sort the smaller and larger parts.
\end{enumerate}

A deterministic implementation may always choose the first or last element as pivot. Unfortunately, this can lead to very poor performance on adversarial inputs. For example, if the array is already sorted and we always choose the first element as pivot, the recursion becomes extremely unbalanced:

\[
(n-1)+(n-2)+\dots+1
=
\Theta(n^2)
\]

Thus deterministic Quick Sort has worst case complexity

\[
T(n)=\Theta(n^2)
\]

The standard randomised version instead chooses the pivot uniformly at random.

\begin{verbatim}
uniform_int_distribution<int>uid(0,a.size()-1);
int pivot=uid(rng);
swap(a[pivot],a[r]);
\end{verbatim}

This tiny modification dramatically changes the behaviour of the algorithm. Intuitively, it becomes very unlikely that we repeatedly choose extremely bad pivots.

\subsubsection*{Runtime Analysis}

Let \(T(n)\) denote the expected runtime of randomised Quick Sort on an array of size \(n\) over all possible inputs.

Suppose the pivot finally ends up at position \(k\), where \(0\le k\le n-1\). Then the recursive subproblems have sizes \(k\) and \(n-1-k\). Since partitioning itself takes n operations, we obtain the recurrence

\[
T(n)
=
T(k)+T(n-1-k)+n
\]

Because the pivot is chosen uniformly at random, every value of \(k\) is equally likely. Taking expectation over all random choices made by the algorithm gives

\[
T(n)
=
n
+
\frac{1}{n}\sum_{k=0}^{n-1}
\Big(
T(k)+T(n-1-k)
\Big)
\]

By sumYmmetry,

\[
\sum_{k=0}^{n-1}T(k)
=
\sum_{k=0}^{n-1}T(n-1-k)
\]

therefore

\[
T(n)
=
\Theta(n)
+
\frac{2}{n}\sum_{k=0}^{n-1}T(k)
\]


\[
S(n)=\sum_{k=0}^{n}T(k)
\]

Then

\[
T(n)=S(n)-S(n-1)
\]

and the recurrence becomes

\[
S(n)-S(n-1)
=
cn+\frac{2}{n}S(n-1)
\]
\[
S(n)
=
\left(1+\frac{2}{n}\right)S(n-1)+cn
\]

\[
\frac{S(n)}{(n+1)(n+2)}
=
\frac{S(n-1)}{n(n+1)}
+
\frac{cn}{(n+1)(n+2)}
\]

Summing this recurrence telescopically gives

\[
\frac{S(n)}{(n+1)(n+2)}
=
\Theta\left(
\sum_{k=1}^{n}\frac{1}{k}
\right)
=
\Theta(\log n)
\]

Therefore,

\[
S(n)=\Theta(n^2\log n)
\]

Finally,

\[
T(n)
=
S(n)-S(n-1)
=
\Theta(n\log n)
\]

Hence the expected runtime of randomised Quick Sort is

\[
T(n)=\Theta(n\log n)
\]

This analysis demonstrated a powerful technique to find the expected runtime - recursion on input size. It was a bit messumY and lengthy, but sometimes yields results very quickly. 

 Another possible way is through indicator variables. Observe that any pair of elements is compared at most once during the entire execution of Quick Sort.

Let

\[
I_{ij}=
\begin{cases}
1,&\text{if elements }i,j\text{ are compared}\\
0,&\text{otherwise}
\end{cases}
\]

Then the total number of comparisons is

\[
C=\sum_{1\le i<j\le n}I_{ij}
\]

Two elements are compared only if one of them becomes the first pivot chosen among all elements lying between them in sorted order. It can be shown that

\[
\Pr(I_{ij}=1)=\frac{2}{j-i+1}
\]

Therefore,

\[
\mathbb{E}[C]
=
\sum_{1\le i<j\le n}\frac{2}{j-i+1}
=
\Theta(n\log n)
\]

Hence randomised Quick Sort has expected complexity

\[
T(n)=\Theta(n\log n)
\]

\subsection{Quick Select}

The Problem: find the \(k\)-th smallest element of an unsorted input array\\

\noindent An obvious method is to sort the array using Merge Sort in $\theta(nlogn)$ and just access the \(k\)-th element - this is overall $\theta(nlogn)$ and is deterministic. Then how could randomness help?\\


Quick Select is an algorithm used to find the \(k\)-th smallest element of an array.

Similar to Quick Sort, the algorithm chooses a pivot and partitions the array into three categories by doing n-1 comparisons between the pivot and every other element as before:
\begin{itemize}
    \item elements smaller than the pivot,
    \item elements equal to the pivot,
    \item elements greater than the pivot.
\end{itemize}

However, unlike Quick Sort, we recurse into only \textbf{one} side:
\begin{itemize}
    \item if the pivot ends up at position \(k\), we are done,
    \item if \(k\) lies to the left, recurse left,
    \item otherwise recurse right.
\end{itemize}

Thus each step discards a significant fraction of the array. This is a similar approach to Binary Search.

\subsubsection*{Expected Runtime}

Suppose partitioning takes n time where array size is n. Let T(n) be the runtime. If
the pivot is the $i$-th smallest element, then 
\[
T(n)=n+
\]
which is a geometric series:

\[
T(n)=\Theta(n)
\]

Thus randomised Quick Select runs in expected linear time.

\section{My Examples}
\subsection{Matrix Problem}

Problem\footnote{This question was taken from Codeforces 1980E} Statement: Given two \(n\times m\) matrices \(A\) and \(B\), determine if they are equivalent under row and column permutations. That is, can we permute the rows and columns of \(A\) to obtain \(B\)?
\\ Additional Constraints: A and B both satisfy -- the set of all entries is \{1,2, \dots, nm\}. \\

\noindent Solution: First we reduce the problem to a simpler one. If we consider R(M) to be the set of sets of values in the rows of a matrix M, and similarly C(M) then the matrices A and B are equivalent if and only if R(A) = R(B) and C(A) = C(B).
This reduced problem can be solved by two approaches:
\subsubsection[Deterministic Method]{Method 1: Set of Sets Approach (Row and Column sets)}

Formally, let
\[
R_i = \{a_{i1}, a_{i2}, \dots, a_{im}\}
\]
be the representation of row \(i\). We compare:
\[
\{R_1, R_2, \dots, R_n\}_A = \{R_1, R_2, \dots, R_n\}_B
\]
and similarly for columns.

This method is conceptually exact but computationally heavy due to the need to compare complex set-of-sets structures.
Set structures can be implemented using balanced binary trees in $\Theta(n\log n)$ time - In \texttt{C++} we have the \texttt{set} data structure which is implemented as a balanced binary tree. 
Comparing two sets of size \(m\) takes \(\Theta(m\log m)\) time, and we have to do this for \(n\) rows and \(n\) columns, leading to a total time complexity of \(\Theta(nm\log n\log m)\). \\

\noindent \texttt{C++} implementation:
\begin{verbatim}
//Matrices are 0-indexed and stored in 2D vectors "first" and "second"
set<set<int>>F;//set of sets for first matrix
for(int i=0;i<n;i++)
{
    set<int>temp(first[i].begin(),first[i].end()); //set of ith row
    F.insert(temp);
}

for(int i=0;i<n;i++)
{
    set<int>temp(second[i].begin(),second[i].end()); //set of ith row
    if (!F.count(temp)) //early exit if this row doesn't exist in the first matrix
    {
        cout<<"NO\n";
        return;
    }
}
\end{verbatim}

\subsubsection[Randomised Method]{Method 2: Hashing Based Verification of Rows and Columns}

To obtain a more efficient solution, each row is hashed to a pair of values $(\text{sum}, \text{squaresum})$, where $\text{sum}$ and $\text{squaresum}$ denote the sum and sum of squares of elements in the row respectively. This representation is much easier to compare across rows.

Assuming the hash is sufficiently collision-resistant, we compare the sets of row hashes for both matrices. This can be implemented efficiently using \texttt{unordered\_set} in \texttt{C++}, which provides expected $O(1)$ insertion and lookup time.

For each row $i$, we compute:
\[
h(R_i) = \left(\sum_{j=1}^{m} a_{ij}, \sum_{j=1}^{m} a_{ij}^2\right)
\]

We define a custom hash for a pair $(x,y)$ as:
\[
H(x,y) = x \cdot 2^{32} + y
\]
where $x$ and $y$ are the computed row statistics. This choice of pair hash is popular for 32 bit integers, as it combines the two values into a single 64-bit integer guaranteeing unique hashes.

We insert the hashes of all rows of matrix \(A\) into an unordered set and then check if the hashes of rows of matrix \(B\) are present in this set. If any hash from \(B\) is not found in the set from \(A\), we can conclude that the matrices are not equivalent under row permutations.

Similarly for columns, we compute the column hashes and compare them in the same manner.

This method has an expected time complexity of \(\Theta(nm)\) due to the linear time required to compute the hashes and the expected constant time for hash set operations.\\

This method involved hashing so the correctness of the algorithm is itself a random variable -- but given the constraint that \\ 
the matrix entries are from \{1,2,\dots,nm\}, it seems highly unlikely that this will fail. \\
In fact, it seems that this may be a unique hash in this case, but i am yet to find a proof for general n,m
\noindent \texttt{C++} implementation:
\begin{verbatim}
auto pair_hash=[](int x,int y){
    return ((x)<<32)+y;
};

unordered_set<long long>s;
loop(i,n)
{
    long long v1=0,v2=0;
    loop(j,m)
    {
        int x=first[i][j];
        v1+=x;
        v2+=x*x;
    }
    s.insert(pair_hash(v1,v2));
}
loop(i,n)
{
    long long v1=0,v2=0;
    loop(j,m)
    {
        int x=second[i][j];
        v1+=x;
        v2+=x*x;
    }
    if(!s.count(pair_hash(v1,v2))) //early exit if this row doesn't exist in the first matrix
    {
        cout<<"NO\n";
        return;
    }          
}
\end{verbatim}

\subsection{Rolling Window Multiset}

\textbf{Problem Statement:} Two integer arrays $x$ and $y$, each of size $n$, are given along with an integer $k$ $(1 \le k \le n)$. For every contiguous subarray (window) of size $k$, determine whether the corresponding windows in $x$ and $y$ contain the same multiset of elements at every position.

More precisely, for every index (1-based) $i$ such that $1 \le i \le n-k+1$, we compare the multisets:
\[
\{x_i, x_{i+1}, \dots, x_{i+k-1}\} \quad \text{and} \quad \{y_i, y_{i+1}, \dots, y_{i+k-1}\}
\]
and check whether they are identical for all $i$.

\subsubsection*{Rolling Hash Method}

To obtain an efficient solution, we maintain two rolling hashes for each window:
\[
S = \sum a_i, \quad Q = \sum a_i^2
\]

These values can be updated in $\mathcal{O}(1)$ time when the window slides by removing the outgoing element and adding the incoming element.

For each position $i$, we compute:
\[
(S_x(i), Q_x(i)) \quad \text{and} \quad (S_y(i), Q_y(i))
\]
and check whether:
\[
(S_x(i), Q_x(i)) = (S_y(i), Q_y(i))
\]

If the condition holds for all $i$, we conclude that every corresponding window has identical multiset structure.\\

\noindent\textbf{C++ Implementation}

\begin{verbatim}
long long sumX=0,sumY=0;
long long sumX2=0,sumY2=0;

for(int i=0;i<k;i++)
{
    sumX+=x[i];
    sumY+=y[i];
    sumX2+=x[i]*x[i];
    sumY2+=y[i]*y[i];
}

bool ok=1;
for(int i=k;i<n;i++)
{
    if(!((sumX==sumY)&&(sumX2==sumY2)))
    {
        ok=0;
        break;
    }

    sumX-=x[i-k];
    sumX+=x[i];
    sumX2-=x[i-k]*x[i-k];
    sumX2+=x[i]*x[i];

    sumY-=y[i-k];
    sumY+=y[i];
    sumY2-=y[i-k]*y[i-k];
    sumY2+=y[i]*y[i];
}
cout<<(ok?"YES":"NO");
\end{verbatim}

This method is very easy to implement and runs in $\mathcal{O}(n)$ time, which is efficient for this problem. \\
However, it relies on the assumption that the hash values are unique for different multisets.\\
We can improve this (if needed) by doing a second round of hashing with a different hash function like polynomial hashing. \\

\end{document}